\section{Formalization design and reproducibility}\label{sec:design}

\subsection{Relation versus division function}

Mathlib already provides polynomial division with respect to a monomial
order.  The relation $\red{G}$ has a different role: it specifies which
single cancellations are mathematically valid, independently of how an
algorithm chooses them.  The current development proves bridge lemmas showing
that a finite relational reduction sequence yields a finite linear-combination
representation of the input, together with the same product-degree bound that
appears in \lean{MonomialOrder.div_set}.  It does not yet prove that the remainder
returned by Mathlib's division procedure is reachable by $\redstar{G}$ and is
normal modulo $G$; establishing that direct connection remains future work.
The relational specification can also be used to compare leading-term,
arbitrary-term, matrix-based, and signature-restricted reduction procedures.

\subsection{Existing and new infrastructure}

The modified \lean{Relation.lean} file contains both pre-existing Mathlib
code and new declarations.  The new declarations include \lean{Diamond},
\lean{Confluent}, \lean{ChurchRosser}, the closure bridges
\lean{Relation.ReflTransGen.to_eqvGen} and \lean{Relation.Join.to_eqvGen}, and
the proofs relating confluence and Church--Rosser.  The file \lean{NormalForm.lean} is separate and defines the
normal-form API used by the polynomial files.  Accordingly, the paper does not
count the complete length of \lean{Relation.lean} as newly formalized code.

\subsection{General coefficient rings}

The coefficient assumptions are layered according to the result being proved.
\begin{itemize}[leftmargin=2em]
  \item Over an arbitrary commutative ring, the development defines reduction,
    characterizes reducibility and normality, proves termination and degree
    bounds, and extracts finite linear-combination certificates such as
    \leandoc{Mathlib/RingTheory/MvPolynomial/PolynomialReductions}{MonomialOrder.exists_linearCombination_of_reflTransGen_with_degree_bound}{exists_linearCombination_of_reflTransGen_with_degree_bound}.
  \item To prove that an arbitrary polynomial multiple of a reducer reduces to
    zero, it is enough for the reducer's leading coefficient to be regular,
    that is, a non-zero-divisor.
  \item The translation lemma, the equivalence between reduction congruence and
    ideal congruence, and the principal confluence/Gr\"obner equivalences use
    the stronger unit-leading-coefficient hypothesis
    \eqref{eq:unit-condition}.
\end{itemize}
Over a coefficient field, Condition~\eqref{eq:unit-condition} is automatic.

\subsection{Reproducibility of the formal artifact}
\label{sec:reproducibility}

The six reviewed Lean files contain no occurrences of \lean{sorry} or
\lean{admit}.  Two immutable snapshots identify the reviewed development:
\begin{itemize}[leftmargin=2em]
  \item the canonical Lean source is commit
    \href{https://github.com/Sanghyeok0/mathlib4/tree/4868fec773dcba035d283904f8b118511cc83d73}
    {\lean{4868fec773dcba035d283904f8b118511cc83d73}}
    of the author's \lean{mathlib4} fork;
  \item the companion artifact is commit
    \href{https://github.com/Sanghyeok0/Buchberger/commit/3e7bbc935df0f545f4d56ed4a457fc23510e30e7}
    {\lean{3e7bbc935df0f545f4d56ed4a457fc23510e30e7}}
    of the \lean{Buchberger} repository.
\end{itemize}
The artifact pins the source commit and the Lean
\lean{v4.33.0-rc1} toolchain.  A clean reproduction uses
\begin{verbatim}
git clone https://github.com/Sanghyeok0/Buchberger.git
cd Buchberger
git checkout 3e7bbc935df0f545f4d56ed4a457fc23510e30e7
lake build
\end{verbatim}
The root module built by this command contains no copied formalization
code: it imports exactly the six modules shown in
Figure~\ref{fig:dependencies} from the pinned fork.  For a direct check of the
final module in the canonical source repository, run
\begin{verbatim}
git clone https://github.com/Sanghyeok0/mathlib4.git
cd mathlib4
git checkout 4868fec773dcba035d283904f8b118511cc83d73
lake exe cache get
lake env lean Mathlib/RingTheory/MvPolynomial/BuchbergerAlgorithm.lean
\end{verbatim}
The artifact homepage
\url{https://sanghyeok0.github.io/Buchberger/} provides reproduction
instructions.  Generated API documentation is available at
\url{https://sanghyeok0.github.io/Buchberger/docs/}; each documented
declaration links to its definition in the fixed \lean{mathlib4} source.
